\subsection{Character values on roots of the full twist}
\label{Section: Character values on roots of the full twist}

For a character $\chi_{\bm v}$ of the generic Hecke algebra $\mathcal{H}_{K(\bm v)}(W)$, let $m^{\chi_{\bm v}}_{\mathcal{C},j}$ denote the multiplicity of $u_{\mathcal{C},j}$ as an eigenvalue of any braid reflection $\bm s_{\mathcal{C},\gamma}$ in the representation $U$ associated with $\chi_{\bm v}$. After Tits' deformation theorem (in particular, after \eqref{EQ: spectra are respected by tits def.}) this equals the multiplicity of $\zeta_{e_{\mathcal{C}}}^j=\sigma (u_{\mathcal{C},j})$ as an eigenvalue of any {\it distinguished} reflection $s_H,\ H\in\mathcal{C}$, in the representation $\op{d}_{\sigma}(U)$. 

The same is true for any admissible specialization $\theta$ (notice that since $f\circ\theta (u_{\mathcal{C},j})=\zeta_{e_{\mathcal{C}}}^j$, the elements $\theta(u_{\mathcal{C},j})$ cannot be equal), so for the analogously defined numbers $m^{\chi_{\bm y}}_{\mathcal{C},j},\ m^{\chi_y}_{\mathcal{C},j},\ m^{\chi}_{\mathcal{C},j}$, we have $$  m^{\chi_{\bm v}}_{\mathcal{C},j}\ =\ m^{\chi_{\bm y}}_{\mathcal{C},j}\ =\ m^{\chi_y}_{\mathcal{C},j}\ =\ m^{\chi}_{\mathcal{C},j}.$$ In view of this, we will only use the latter notation $m^{\chi}_{\mathcal{C},j}$ from now on. Notice finally that by the defining relations \eqref{EQ: deformed order relations}, the only possible eigenvalues for any $\bm s_{\mathcal{C},\gamma}$ are precisely the $u_{\mathcal{C},j}$'s. We therefore have (for any $\mathcal{C}\in\mathcal{A}/W$) \begin{equation} \sum_{j=0}^{e_{\mathcal{C}}-1}m^{\chi}_{\mathcal{C},j}=\chi(1).\label{Eq: sum of m_H,j's=chi(1)}\end{equation} 


The following proposition is essential for the proof of our technical lemma (Prop.~\ref{Prop: The key technical lemma}). To simplify its statement we first introduce the following notation (recall also that $\omega_{\mathcal{C}}=|\mathcal{C}|$ for an orbit $\mathcal{C}\in\mathcal{A}/W$):

\begin{defi}\label{Def. Pre character value stuff}
Consider\footnote{We move to a larger ring, so that the roots $u_{\mathcal{C},j}^{1/\chi(1)}$ are well defined. Shortly however, Prop.~\ref{Prop: character values on roots of full twist} will show that these monomials actually live in $K[\bm v]$.} the element of $K[\bm u^{1/|W|}]$ given as $$z_{\chi_{\bm v}}(\bm \pi)\coloneqq \prod_{\mathcal{C}\in\mathcal{A}/W}\prod_{j=0}^{e_{\mathcal{C}}-1}u_{\mathcal{C},j}^{(1/\chi(1))m^{\chi}_{\mathcal{C},j}e_{\mathcal{C}}\omega_{\mathcal{C}}},$$ and, for a regular number $d$ \textup{(}see Defn.~\ref{Defn: regular element}\textup{)}, write $$z_{\chi_{\bm v}}(\bm \pi)^{1/d}\coloneqq \prod_{\mathcal{C}\in\mathcal{A}/W}\prod_{j=0}^{e_{\mathcal{C}}-1}u_{\mathcal{C},j}^{(1/d\chi(1))m^{\chi}_{\mathcal{C},j}e_{\mathcal{C}}\omega_{\mathcal{C}}}.$$
Finally, denote by $N(\chi)$ the quantity $$N(\chi)\coloneqq \sum_{\mathcal{C}\in\mathcal{A}/W}\omega_{\mathcal{C}}\cdot\sum_{j=0}^{e_{\mathcal{C}}-1}jm^{\chi}_{\mathcal{C},j}.$$
\end{defi}

\begin{rema}\label{Remark: N(x) via coexponents}
$N(\chi)$ usually denotes the sum of the $\chi^{*}$-exponents (see \cite[Chapter~4:~\S4]{lehrer-taylor-book}) of the representation that affords $\chi$. This in fact agrees with the definition above (see \cite[Prop.~4.1]{broue-michel-french-paper}, or \cite[Lemma~10.15 and Remark~10.12]{lehrer-taylor-book} which includes Gutkin's theorem). We are only going to use it as a symbol (but see also Remark~\ref{Remark: cox numbers and N(x)+N(x*)}).
\end{rema}

As we have mentioned earlier, the reason that we have nice, explicit character evaluations on the full twist $T_{\bm\pi}$ and its roots is that $\bm\pi$ is central in $B(W)$. Take for instance the determinant character $\op{det}_{\chi_{\bm v}}$ associated to $\chi_{\bm v}$. It is linear and therefore factors through the abelianization $B^{\op{ab}}$ so that Corol.~\ref{Cor: the abelianization of the full twist} implies that its values on powers of the full twist are given by 
\[
\op{det}_{\chi_{\bm v}}(T_{\bm \pi}^l)=\prod_{\mathcal{C}\in\mathcal{A}/W}\prod_{j=0}^{e_{\mathcal{C}}-1}u_{\mathcal{C},j}^{m^{\chi}_{\mathcal{C},j}\omega_{\mathcal{C}}e_{\mathcal{C}}l}=z_{\chi_{\bm v}}(\bm \pi)^{\chi(1)l}.
\]

Now, since $T_{\bm\pi}^l$ is central, it acts on irreducible representations as a scalar. That is, its spectrum is given by $$\op{Spec}_{\chi_{\bm v}}(T_{\bm\pi}^l)=\{\xi z_{\chi_{\bm v}}(\bm\pi)^l\ (\chi(1)\text{-many times})\},$$ where $\xi$ is a $\chi(1)$-th root of unity. This works similarly for roots of $T_{\bm\pi}$ and with Tits' deformation theorem and little more work we get the following.

\begin{prop}[{\cite[Prop.~4.16]{broue-michel-french-paper}}]\label{Prop: character values on roots of full twist}
For a character $\chi_{\bm v}$ of the generic Hecke algebra, the values on the full twist $T_{\bm\pi}$ are given by $$ \chi_{\bm v}(T_{\bm\pi})=\chi(1)e^{-2i\pi N(\chi)/\chi(1)}z_{\chi_{\bm v}}(\bm \pi).$$
Moreover, if $\bm w$ is a $d$-th root of some power $\bm\pi^l$ and its image in $W$ under the fixed surjection \eqref{1PBW1} is $w$, we have $$\chi_{\bm v}(T_{\bm w})=\chi(w)e^{-2i\pi lN(\chi)/d\chi(1)}z_{\chi_{\bm v}}(\bm\pi)^{l/d}. $$
\end{prop}

\begin{rema}
In \cite{broue-michel-french-paper} the Hecke algebras are introduced \emph{formally} by deforming the Artin-like presentations of the generalized braid groups, while we have used the topological interpretation of \cite{broue-malle-rouquier-braid-hecke}. This does not affect the proof of the previous proposition which only relies on the centrality of the full twist $\bm\pi$ and Corol.~\ref{Cor: the abelianization of the full twist}.
\end{rema}


By applying the specialization $\theta_{\bm x}$ from \eqref{EQ: admissible specializations} to the previous proposition, we easily get: 

\begin{coro}\label{Corol: character calculation for H_x(W)}
Let $\bm w$ be a $d$-th root of some power $\bm \pi^l$ as above and let $\chi_{\bm y}$ be a character of the specialization $\mathcal{H}_{\bm x}(W)$ as in \eqref{EQ: admissible specializations and Tits factoring}. We have \begin{enumerate}
\item\ \vspace{-0.35cm} $$\chi_{\bm y}(T_{\bm \pi})=\chi(1)\prod_{\mathcal{C}\in\mathcal{A}/W}x_{\mathcal{C}}^{(1/\chi(1))m^{\chi}_{\mathcal{C},0}e_{\mathcal{C}}\omega_{\mathcal{C}}}.$$
\item\ \vspace{-0.35cm} $$\chi_{\bm y}(T_{\bm w})=\chi(w)\prod_{\mathcal{C}\in\mathcal{A}/W}x_{\mathcal{C}}^{(l/d\chi(1))m^{\chi}_{\mathcal{C},0}e_{\mathcal{C}}\omega_{\mathcal{C}}}.$$
\end{enumerate}
\end{coro}


\subsection{Local Coxeter numbers}

We are now going to define a local version of Coxeter numbers (see Defn.~\ref{Defn: Coxeter numbers}) and study how they are precisely related to the exponents that appear in the character calculation of the previous Corol.~\ref{Corol: character calculation for H_x(W)}.


\begin{defi}\label{Defn: Local Coxeter numbers} 
We define the \emph{local} Coxeter number $c_{\chi_{,\mathcal{C}}}$ associated to the character $\chi$ and the hyperplane orbit $\mathcal{C}\in\mathcal{A}/W$, as the normalized trace $$c_{\chi_{,\mathcal{C}}}\coloneqq \dfrac{1}{\chi(1)}\cdot\chi\Big( \sum_{V^{t}\in\mathcal{C}}(\bm 1-t)\Big).$$ Here, the sum is taken over all reflections $t$ whose fixed hyperplane $H=V^{t}$ belongs to the orbit $\mathcal{C}$. Notice that these numbers are a refinement of the Coxeter numbers in the sense that $c_{\chi}=\sum c_{\chi_{,\mathcal{C}}}$
\end{defi}

\begin{prop}\label{Prop: local cox num formula}
The local Coxeter numbers satisfy $$c_{\chi_{,\mathcal{C}}}=e_{\mathcal{C}}\cdot\omega_{\mathcal{C}}\cdot\Big(1-\dfrac{m^{\chi}_{\mathcal{C},0}}{\chi(1)}\Big).$$
\end{prop}
\begin{proof}
As we saw in \eqref{EQ: Reflection set partition into cyclis}, because the parabolic groups for hyperplanes are cyclic, the set of reflections can be partitioned into sets of the form $\{t_H,\ldots,t^{e_H-1}_H\}$. Moreover, recalling the definition of $m^{\chi}_{\mathcal{C},j}$ from the beginning of this section, we see that the spectrum of $t_H^k$ (for $H\in\mathcal{C}$) is given by $$\op{Spec}_{\chi}(t_H^k)=\{ \zeta_{e_{\mathcal{C}}}^{jk}\ (m^{\chi}_{\mathcal{C},j}\text{-many times})\ |\ 0\leq j\leq e_{\mathcal{C}}-1\}.$$ We can then pick an $H\in\mathcal{C}$ and a generator $t_H$ of $W_H$, and start the evaluation by computing \begin{align*}
\sum_{V^t\in\mathcal{C}}\chi(\bm 1-t)&=\chi(1)(e_{\mathcal{C}}-1)\omega_{\mathcal{C}}-\omega_{\mathcal{C}}\sum_{k=1}^{e_{\mathcal{C}}-1}\chi(t_H^k)\\
&=\chi(1)(e_{\mathcal{C}}-1)\omega_{\mathcal{C}}-\omega_{\mathcal{C}}\sum_{k=1}^{e_{\mathcal{C}}-1}\sum_{j=0}^{e_{\mathcal{C}}-1}m^{\chi}_{\mathcal{C},j}\zeta_{e_{\mathcal{C}}}^{jk}.
\end{align*}
Now, notice that the sum $\sum_{k=1}^{e_{\mathcal{C}}-1}\zeta_{e_{\mathcal{C}}}^{jk}$ equals $e_{\mathcal{C}}-1$ or $-1$ depending on whether $j=0$ or not. So, after changing the order of summation, we have \begin{align*}
\sum_{V^t\in\mathcal{C}}\chi(\bm 1-t)&=\chi(1)(e_{\mathcal{C}}-1)\omega_{\mathcal{C}}+\sum_{j=1}^{e_{\mathcal{C}}-1}\omega_{\mathcal{C}}m^{\chi}_{\mathcal{C},j}-(e_{\mathcal{C}}-1)\omega_{\mathcal{C}}m^{\chi}_{\mathcal{C},0}\\
&=\chi(1)e_{\mathcal{C}}\omega_{\mathcal{C}}-e_{\mathcal{C}}\omega_{\mathcal{C}}m^{\chi}_{\mathcal{C},0},
\end{align*}where the second equality is because of \eqref{Eq: sum of m_H,j's=chi(1)}. This completes the proof.
\end{proof}

We can now rewrite the character calculation from Corol.~\ref{Corol: character calculation for H_x(W)} replacing the quantities in the exponents with equivalent ones in terms of the Coxeter numbers $c_{\chi_{,\mathcal{C}}}$ (and via Prop.~\ref{Prop: local cox num formula}). With the notation being the same as in the statement of the Corollary, we have: \begin{equation}
\chi_{\bm y}(T_{\bm \pi})=\chi(1)\prod_{\mathcal{C}\in\mathcal{A}/W}x_{\mathcal{C}}^{e_{\mathcal{C}}\omega_{\mathcal{C}}-c_{\chi_{,\mathcal{C}}}}\hspace{1.5mm}\text{ and }\hspace{1.5mm}\chi_{\bm y}(T_{\bm w})=\chi(w)\prod_{\mathcal{C}\in\mathcal{A}/W}x_{\mathcal{C}}^{(e_{\mathcal{C}}\omega_{\mathcal{C}}-c_{\chi_{,\mathcal{C}}})l/d}.\label{EQ: character values via local Coxeter}\end{equation}

Moreover, after the further specialization $x_{\mathcal{C}}\mapsto x$ of $\theta_x$ from \eqref{EQ: admissible specializations} and for the characters $\chi_y$ of $\mathcal{H}_x(W)$ as in \eqref{EQ: admissible specializations and Tits factoring}, we have (recalling that $\sum e_{\mathcal{C}}\omega_{\mathcal{C}}=|\mathcal{R}|+|\mathcal{A}|$ and that $\sum c_{\chi_{,\mathcal{C}}}=c_{\chi}$):
\begin{equation}
\chi_y(T_{\bm \pi})=\chi(1)\cdot x^{|\mathcal{R}|+|\mathcal{A}|-c_{\chi}}\quad\text{ and }\quad\chi_y(T_{\bm w})=\chi(w)\cdot x^{(|\mathcal{R}|+|\mathcal{A}|-c_{\chi})l/d}.\label{EQ: character values via Coxeter-not local}\end{equation}


\begin{rema}\label{Remark: cox numbers and N(x)+N(x*)}
This last equation is precisely what appears in \cite[Prop.~4.18]{broue-michel-french-paper} but with an equivalent expression for the Coxeter numbers: $$c_{\chi}=\dfrac{N(\chi)+N(\chi^*)}{\chi(1)},$$ where the numbers $N(\chi)$ are given in Defn.~\ref{Def. Pre character value stuff} (see also Rem.~\ref{Remark: N(x) via coexponents}). This expression also appears in \cite[Lemma~1]{michel-deligne-lusztig-derivation} but the statement of that lemma might be misleading as it holds regardless of the values $e_{\mathcal{C}}$. For completeness, we include the calculation: $$\chi(1)c_{\chi}=\chi(1)\sum_{\mathcal{C}\in\mathcal{A}/W}c_{\chi_{,\mathcal{C}}}=\sum_{\mathcal{C}\in\mathcal{A}/W}\omega_{\mathcal{C}}\sum_{j=1}^{e_{\mathcal{C}}-1}e_{\mathcal{C}}m^{\chi}_{\mathcal{C},j}=N(\chi)+N(\chi^*).$$
In fact, Michel later on \cite[Rem.~2]{michel-deligne-lusztig-derivation} notes that for all groups $W$ one has (see Defn.~\ref{Defn: Malle's character permutations} for Malle's permutation $\Psi$ on the set of irreducible characters $\chi$ of $W$) $$c_{\chi}=\dfrac{N(\chi)+N(\Psi(\chi^*))}{\chi(1)},$$ which is equivalent to the first statement as $N(\Psi(\chi))=N(\chi)$ after Prop.~\ref{Prop: properties of Psi for Cox numbers}.
\end{rema}

The following generalizes Rem.~\ref{Rem: bounds Coxeter nums} and is a direct corollary of Prop.~\ref{Prop: local cox num formula}.

\begin{coro}\label{Corol: local cox num are integers andbounds}
The Coxeter numbers $c_{\chi_{,\mathcal{C}}}$ are integers and they satisfy $$ 0\leq c_{\chi_{,\mathcal{C}}}\leq e_{\mathcal{C}}\cdot\omega_{\mathcal{C}}.$$
\end{coro}
\begin{proof}
The inequalities are immediate from Prop.~\ref{Prop: local cox num formula}, since $0\leq m^{\chi}_{\mathcal{C},0}\leq \chi(1)$. For the integrality property, notice first that the collection of reflections $t\in \mathcal{R}$ such that $V^t\in\mathcal{C}$, is a union of conjugacy classes, so that the numbers $c_{\chi_{,\mathcal{C}}}$ are in fact \emph{algebraic integers} \cite[Corol.~1, p.~52]{Serre}. After Prop.~\ref{Prop: local cox num formula} they are also clearly rational numbers and the result follows.
\end{proof}

